\section{独立练习}

\begin{question}\begin{drill}[弱偏转的最近距离修正]\label{at4:drill:1}
入射参数 $b,v_\infty$ 满足 $\eta=\mu/(bv_\infty^2)\ll1$，求偏转角和最近距离的一阶近似。
\end{drill}
\typeout{ASTROFIG:fig:at4:drill:1:s}\fig{fig:at4:drill:1:s}{题设：弱偏转的最近距离修正}{\begin{tikzpicture}[scale=.85]
\draw[guide](-4.5,-1.1055)--(4.5,-1.1055);\draw[world,domain=-2:2,samples=100]plot({(.11111111*(10-cosh(\x))+9.94987437*1.1055416*sinh(\x))/10},{(-9.94987437*.11111111*(10-cosh(\x))+1.1055416*sinh(\x))/10});\fill(0,0)circle(2pt)node[above]{$F$};\draw[flow](-4.1,-1.1055)--(-3,-1.1055)node[below]{$v_\infty$};
\end{tikzpicture}}
\end{question}
\hintlink{at4:drill:1}
\begin{question}\begin{drill}[给定余速与安全半径的最大转角]\label{at4:drill:2}
行星参数 $\mu_p$，相对余速 $u$，掠近点不得小于 $r_{\min}$。求最大偏转角和对应冲量参数，并说明能否精确反转相对速度。
\end{drill}
\typeout{ASTROFIG:fig:at4:drill:2:s}\fig{fig:at4:drill:2:s}{题设：给定余速与安全半径的最大转角}{\begin{tikzpicture}[scale=.85]
\draw[guide](-4.6,-1.118)--(1,-1.118);\draw[world,domain=-2:2,samples=110]plot({(1.5-cosh(\x))/1.5+1.11803399*1.11803399*sinh(\x)/1.5},{-1.11803399*(1.5-cosh(\x))/1.5+1.11803399*sinh(\x)/1.5});\fill(0,0)circle(2pt)node[above left]{$F$};\draw[flow](-4.2,-1.118)--(-3,-1.118)node[midway,below]{$u$};\draw[world](0,0)circle(.5);
\end{tikzpicture}}
\end{question}
\hintlink{at4:drill:2}
\begin{question}\begin{drill}[引力助推的能量正负]\label{at4:drill:3}
行星在恒星系速度为 $\bm V$。探测器入、出行星系余速大小均为 $u$，分别与 $\bm V$ 夹角 $\alpha,\beta$。求恒星系单位质量能量变化，并给出加速条件。
\end{drill}
\typeout{ASTROFIG:fig:at4:drill:3:s}\fig{fig:at4:drill:3:s}{题设：引力助推的能量正负}{\begin{tikzpicture}[scale=.85]
\draw[flow](0,0)--(2,0)node[below]{$\bm V$};\draw[flow](0,0)--(1,1.5)node[above]{$\bm u_{\rm in}$};\draw[flow,blue](0,0)--(1.7,.5)node[right]{$\bm u_{\rm out}$};
\end{tikzpicture}}
\end{question}
\hintlink{at4:drill:3}
\begin{question}\begin{drill}[局部速度相同却有不同最近距离]\label{at4:drill:4}
两粒子在相同半径 $r_0$，速率均为 $\sqrt{\mu/r_0}$。一个纯切向，另一个径向向外。比较两条轨道的能量、角动量和最大最小距离。
\end{drill}
\typeout{ASTROFIG:fig:at4:drill:4:s}\fig{fig:at4:drill:4:s}{题设：局部速度相同却有不同最近距离}{\begin{tikzpicture}[scale=.85]
\fill(0,0)circle(2pt)node[below]{$F$};\draw[guide](0,0)circle(2);\fill(2,0)circle(2pt);\draw[flow](2,0)--(2,1.3)node[left]{切向};\draw[flow,blue](2,0)--(3.3,0)node[above]{径向};\draw[dim](0,-.6)--(2,-.6)node[midway,below]{$r_0$};
\end{tikzpicture}}
\end{question}
\hintlink{at4:drill:4}
\begin{question}\begin{drill}[无耗散引力能否自动捕获]\label{at4:drill:5}
粒子从无穷远以非零余速掠过一个静止球形天体，轨道不碰表面。仅靠二体保守引力，能否最终留在椭圆轨道？若能量从其他系统交换，应如何修改判断？
\end{drill}
\typeout{ASTROFIG:fig:at4:drill:5:s}\fig{fig:at4:drill:5:s}{题设：无耗散引力能否自动捕获}{\begin{tikzpicture}[scale=.85]
\draw[guide](-4.6,-1.118)--(1,-1.118);\draw[world,domain=-2:2,samples=110]plot({(1.5-cosh(\x))/1.5+1.11803399*1.11803399*sinh(\x)/1.5},{-1.11803399*(1.5-cosh(\x))/1.5+1.11803399*sinh(\x)/1.5});\fill(0,0)circle(2pt)node[above left]{$F$};\draw[flow](-4.2,-1.118)--(-3,-1.118)node[midway,below]{$v_\infty$};
\end{tikzpicture}}
\end{question}
\hintlink{at4:drill:5}
\begin{question}\begin{drill}[超过给定偏转角的散射截面]\label{at4:drill:6}
固定中心参数 $\mu$，入射速度 $u$，统计偏转角至少为 $\chi_0$ 的轨道，$0<\chi_0<\pi$。求点质量模型下的截面；若中心有吸收半径，再说明怎样排除碰撞。
\end{drill}
\typeout{ASTROFIG:fig:at4:drill:6:s}\fig{fig:at4:drill:6:s}{题设：超过给定偏转角的散射截面}{\begin{tikzpicture}[scale=.85]
\draw[guide](-4.6,-1.118)--(1,-1.118);\draw[world,domain=-2:2,samples=110]plot({(1.5-cosh(\x))/1.5+1.11803399*1.11803399*sinh(\x)/1.5},{-1.11803399*(1.5-cosh(\x))/1.5+1.11803399*sinh(\x)/1.5});\fill(0,0)circle(2pt)node[above left]{$F$};\draw[flow](-4.2,-1.118)--(-3,-1.118)node[midway,below]{$u$};
\end{tikzpicture}}
\end{question}
\hintlink{at4:drill:6}
\begin{question}\begin{drill}[吸收静止尘埃后的减速与耗散]\label{at4:drill:7}
质量 $M$ 的天体以速度 $V$ 运动，短时间完全吸收原先静止的质量 $\dd M$。忽略这段时间外力冲量，求一阶速度变化和机械能损失。
\end{drill}
\typeout{ASTROFIG:fig:at4:drill:7:s}\fig{fig:at4:drill:7:s}{题设：吸收静止尘埃后的减速与耗散}{\begin{tikzpicture}[scale=.85]
\draw[world](0,-.4)rectangle(1.5,.4);\node at(.75,0){$M$};\draw[flow](0,1)--(1.5,1)node[above]{$V$};\fill(3,0)circle(2pt)node[above]{$\dd M$ 静止};
\end{tikzpicture}}
\end{question}
\hintlink{at4:drill:7}
\begin{question}\begin{drill}[恒星系助推能量的上界为何不能随意取得]\label{at4:drill:8}
行星日心速度大小 $V$，探测器相对余速 $u$。证明任一次无动力掠近的单位质量日心能量变化满足 $|\Delta\varepsilon|\le2Vu$；若转角最多 $\delta_{\max}$，给出更紧的方向无关上界。
\end{drill}
\typeout{ASTROFIG:fig:at4:drill:8:s}\fig{fig:at4:drill:8:s}{题设：恒星系助推能量的上界为何不能随意取得}{\begin{tikzpicture}[scale=.85]
\draw[guide](0,0)circle(1.6);\draw[flow](0,0)--(1.6,0)node[right]{$\bm u_{\rm in}$};\draw[flow,blue](0,0)--(-.8,1.386)node[above]{$\bm u_{\rm out}$};
\end{tikzpicture}}
\end{question}
\hintlink{at4:drill:8}
